\chapter{The Hard Problem and Its Discontents}
\label{ch:hard_problem}

\papernote{Draws from the introductory analysis in: \emph{On the Consequences of Categorical Completion on Cognitive Processes} and \emph{On the Geometry of Consciousness}}

\begin{chapterabstract}
The hard problem of consciousness asks why physical processes give rise to subjective experience. Three classes of prior solutions exist: eliminativism (consciousness is an illusion), dualism (consciousness is non-physical), and panpsychism (consciousness is a fundamental property). Each fails for structural reasons identified here. A fourth class---functionalism and its descendants, including global workspace theory and integrated information theory---is shown to share a common defect: it specifies the functional correlates of consciousness without explaining why those correlates produce experience at all. This chapter establishes the precise nature of what a successful theory must explain and what it is not permitted to assume.
\end{chapterabstract}

\input{../docs/categorical-completion/sections/hard-problem-survey.tex}
